\newpage
\section{The first-entry polynomial and rank reduction}
Put \(b_{j;m}=u_{j;m}^{-1}\) and \(h_x=[M_e]\mathbf C_x\). In the word convention of the parabolic paper,
the negative, zero and positive shifts in our sorted window give the letters
\(\mathsf C,\mathsf B,\mathsf A\), respectively. Thus the word of \(b_{j;m}\) is
\[
 \mathsf C^{j-1}\mathsf B^{m-j}\mathsf A\mathsf B\mathsf C
                  \mathsf B^{m-j}\mathsf A^{j-1}.
\]
The two reflection chains in that paper keep \(m\) letters of each wound type throughout.
The off-diagonal endpoints have \(j<m\) of each.

\subsection*{The first entry admits a finite Hecke calculation}
\textbf{Deduction.} For every \(m\ge2\),
\[
   P_{w_0,C_{1;m}}(q)=1+q+\cdots+q^{m-1},\qquad \mu_{1;m}=1 .
\]
Indeed, the window of \(s_0\) has \(0\) at position \(1\) and \(N+1\) at
position \(N\). Moving them to positions \(m,m+2\) gives
\[
 b=b_{1;m}=s_0(s_1\cdots s_{m-1})(s_{N-1}\cdots s_{m+2})\in W_J .
\]
Use the same \(V=I\cap J\) and \(L=V\setminus\{s_1,s_{N-1}\}\) as in our gallery.
Thus \(W_V=S_m\times S_m\). Write \([a]=(v^a-v^{-a})/(v-v^{-1})\).
In the spherical right module consider the bar-invariant vector
\(G=M_e C_{s_0}C_V/\pi_L\), where \(C_V=C_{w_V}\).
The right \(W_V\)-orbit of the left coset indexed by \(s_0\) has minimal
index \(s_0\) and maximal index \(b\). Its elements move the one
\(\mathsf A\) within the first block of \(m\) positions and the one
\(\mathsf C\) within the last block, keeping the middle \(\mathsf B\) fixed.
The stabilizer is \(W_L\). Expanding the parabolic symmetrizer into
minimal coset representatives gives
\[
 G=\sum_{x\in\mathcal O_V(s_0)}
 v^{\ell(b)-\ell(x)}M_x+v[m]^2M_e .
\]
Here \(M_e C_{s_0}=M_{s_0}+vM_e\), and \(\pi_V/\pi_L=[m]^2\).
For the orbit term, each shortest shuffle moves the two marked letters
away from their initial positions by length-increasing steps. Thus it
acts by \(M_{s_0}H_z=M_{s_0z}\), with exponent
\(\ell(w_V)-\ell(w_L)-\ell(z)=\ell(b)-\ell(s_0z)\).
This is also the orbit formula underlying B \S3.3.
Now \(v[m]^2-[m][m-1]=v^m[m]\). The vector
\[
 \mathbf C_b=G-[m][m-1]M_e
\]
is characterized by bar invariance, leading standard \(M_b\) with
coefficient one, and strictly positive powers of \(v\) at all other
standards. This proves the displayed KL formula since
\(h_b=v^m[m]\) and \(\ell(b)=2m-1\).

This finite calculation gives the entire first-column KL polynomial. For example, at $m=3$ it gives $1+q+q^2$. Section~\ref{sec:partialrows} determines the leading coefficient of every column.


\subsection*{The first seed in the uncorrected parabolic family}
The same entry can be read from the parabolic product in the ambient period
\(2m+1\). With \(T_1,T_2\) as in \((*)\) and our strong one-line
canonical \(\mathcal E_{1,m}\) refined to \(V\), one has
\[
 \mathcal{BS}_1T_1T_2=\mathbf C_{b_{1;m}}+[m-1]\mathcal E_{1,m}.
\]
Here \(\mathcal{BS}_1=D(J_1)\) has coefficient \(v^{2m-1}\) at \(e\)
and \(v^{c-3}\) at the rank-one words with \(\mathsf A\) in position
one and \(\mathsf C\) in position \(3\le c\le N\) (B Lemma 3.19).
For the three output orbit maxima \(H,U,F\), the \(\mathsf A\) is in
position \(m\) and the \(\mathsf C\) in positions \(1,m+1,m+2\),
respectively. The four standard coordinates are
\[
\begin{array}{c|cccc}
 &H&U&F&e\\ \hline
 \mathcal{BS}_1T_1T_2 &[m-1]&v^{m-1}[m-1]&v^{m-1}[m]&v^{2m-1}[m]^2\\
 \mathcal E_{1,m} &1&v^{m-1}&v^m&v^{2m}[m].
\end{array}
\]
For the first row, \(T_1\) sends the terms \(3\le c\le m+1\) to
their orbit sum with top coefficient
\(\sum_{c=3}^{m+1}v^{c-3+c-m-1}=[m-1]\), and multiplies the others
by \([m]\). For \(H\), the second averaging contributes the sum
\(\sum_{c=2}^m v^{c-2+c-m}=[m-1]\) divided by \([m-1]\);
for \(U,F\) it moves a single \(\mathsf A\) to the end of the first
block. The second row follows from A Lemmas 4F.2 and 4D.1.
These are all occupied orbits. Subtract the rows with the stated factor,
using \(v^{m-1}[m]-v^m[m-1]=1\).

\newpage
\subsection*{The rank-reduction identity}
Let \(\Theta_k\) have word \(\mathsf C^k\mathsf B^{N-2k}\mathsf A^k\).
It is the finite-orbit maximum of type \(k\) (B Lemma 3.11), hence represents
our full-spherical standard \(Z_k\) after output refinement. The coefficient
at its longest fine representative has shift zero (A Lemma 4D.1).
For \(0\le k<j\le m\) we obtain the all-degree identity
\[
   [\Delta_{Z_k}]C_{A_{j;m}}
   =[M_{\Theta_k}]\mathbf C_{b_{j;m}}
   =h^{(m-k)}_{b_{j-k;m-k}},
\]
where the superscript gives the smaller rank parameter.
For \(k>0\), both ordinary endpoints \(w_I\Theta_k,w_Ib_{j;m}\) take
the values \(2N,1-N\) at positions \(1,N\), respectively. Every
fundamental-window entry lies between those bounds. Deleting the common
first and last letters removes one \(\mathsf C\) and one \(\mathsf A\):
the relabelling in B Definition 4.17 gives
\(w'_I\Theta_{k-1},w'_I b_{j-1;m-1}\), with period \(N-2\).
To check Bruhat comparability use the rank functions
\(\#\{i\le a:w(i)\ge b\}\).
At each cut the deleted periodic strands contribute equally to the
two functions. On the remaining strands their inequalities are exactly
the smaller-period inequalities. Iterating reaches the comparable pair
\(w'_I,w'_I b_{j-k;m-k}\). B Proposition 4.18 now applies at each
deletion and preserves the normalized coefficient. For \(j\le k\) the
original coordinate vanishes since
\(\ell(\Theta_k)=2k(N-k)\ge 2j(N-j)>D_{j;m}\).

On the diagonal, the coefficient formula in the parabolic paper gives \([v][\Delta_{Z_k}] C_{A_{m;m}}=m-k\) for
\(0\le k<m\). The final case reduces to \(s_0\) in period three.
The independent KL recursion checked the full polynomial identity for every
\(k,j\) in ranks \(m=2,3,4\). For example, at \(m=4\) the rows of
linear coefficients, with columns \(j=1,2,3,4\), are
\((1,2,3,4),(0,1,2,3),(0,0,1,2),(0,0,0,1)\).
Together with our one-line degree bound at each \(Z_k\) (Lemma 4H.1),
these tests at full-spherical standards recover the row sums in
parameters \(m-k\). To distinguish entries within a fixed row,
Section~\ref{sec:partialrows} computes its successive partial sums.

\subsection*{The length gap in direct rank reduction for \(2\le j<m\)}
At the ordinary pair \(w_I,w_Ib_{j;m}\) the only common window
position is the central \(\mathsf B\) at \(m+1\). The other two
\(\mathsf B\)-blocks have \(b(i)=i+1\) and \(i-1\), respectively.
Deleting the common position and image leaves type \(j\) in period
\(N-1\). In B Lemma 3.9, \(\ell=2j(N-j)-\operatorname{inv}\) with
\(\mathsf C<\mathsf B<\mathsf A\). Deletion drops two inversions, one
with the preceding \(\mathsf A\), one with the following \(\mathsf C\).
The ordinary length gap therefore drops by \(2j-2>0\), preventing
Henderson cancellation at this strand. Across fixed \(j\) the full gap
difference is \(D_{j;m}-D_{j;j}=(4j-2)(m-j)\).

The seed correction also depends on the ambient parameter. B Proposition 3.21 gives
lowest exponent \(2-(N-2j)=1-2(m-j)<0\) for \(\mathcal{BS}_j\) at the
orbit maximum \(z(j,1,1)\). The terms \(\mathcal{BS}_{j-1}\) and
\(\mathcal{CBS}_s\) for \(s\le j-1\) vanish on this type-\(j\) orbit.
Replacing \(m\) by \(j\) in the correction while retaining period \(2m+1\)
therefore leaves a negative power at this orbit, so it does not give the
canonical seed.
